\providecommand{\Ch}{\operatorname{Ch}}
\providecommand{\Chodd}{\operatorname{Ch}_{\mathrm{odd}}}
\providecommand{\Kc}{K_{\mathrm{c}}}
\providecommand{\inda}{\operatorname{ind}_{\mathrm{a}}}
\providecommand{\Torus}{\mathbb{T}}

\section{Solutions for \cref{ch:families}: Multiplicativity, the Chern character and the families index theorem}\label{solutions:ch:families}

\subsection*{\Cref{families:ex:ext-product}}
\emph{Idea: every move generating the equivalence relation of $\Kc$, applied to one
factor, changes $a_1\#a_2$ by a move of the same kind; this uses only
\cref{families:lem:clifford}. The flip changes $a_1\#a_2$ by a sign on one summand.}

Write $m(p,\alpha)=m(1,1)$ formed with $p$ and $\alpha$, so $a_1\#a_2=m(a_1,a_2)$. Let
$K_i$ be compact with $a_i$ invertible on $X_i\setminus K_i$. By
\cref{families:lem:clifford}, $m(p,\alpha)$ is invertible at every point where $p$
\emph{or} $\alpha$ is invertible; this holds for any choice of Hermitian metrics used
to form the adjoints. We check that the class of $\bigl((\mathbb
E^1\otimes\mathbb E^2)_\pm,a_1\#a_2\bigr)$ does not change under the moves applied
to the first factor; the second factor is treated in the same way, the roles of $p$
and $\alpha$ in \cref{families:lem:clifford} being symmetric.

\emph{(a) Extension and changes on compact sets.} If $a_1,a_1'$ are two smooth
extensions which agree outside a compact $K\supset K_1$, then
$a_1^s=(1-s)a_1+sa_1'$ is invertible outside $K$, and $m(a_1^s,a_2)$, $s\in[0,1]$, is
invertible outside $K\times K_2$. So it is a homotopy of triples on $X_1\times X_2$.
This covers both the choice of extension and the move ``change $a$ on a compact
set''.

\emph{(b) Metrics.} If $h^0$, $h^1$ are two Hermitian metrics on the bundles
involved, $h^s=(1-s)h^0+sh^1$ are metrics, the adjoints with respect to $h^s$ depend
smoothly on $s$, and by \cref{families:lem:clifford} (for the metric $h^s$) the
corresponding maps $m^s$ are invertible outside $K_1\times K_2$: a homotopy of triples.
So the class does not depend on the metrics.

\emph{(c) Bundle isomorphisms.} Let $f_\pm:E^1_\pm\to E'_\pm$ be isomorphisms with
$a_1'=f_-a_1f_+^{-1}$ outside a compact set. By (b) we may give $E'_\pm$ the
metrics transported by $f_\pm$, so that $f_\pm$ are unitary and
$(a'_1)^*=f_+a_1^*f_-^{-1}$. Then, with $g_+=f_+\otimes1\oplus f_-\otimes1$ on
$(\mathbb E^1\otimes\mathbb E^2)_+$ and $g_-=f_-\otimes1\oplus f_+\otimes1$ on
$(\mathbb E^1\otimes\mathbb E^2)_-$, each entry of \eqref{families:eq:clifford}
transforms correctly, so $m(a_1',a_2)=g_-\,m(a_1,a_2)\,g_+^{-1}$: the two product
triples are isomorphic.

\emph{(d) Homotopies.} A homotopy of triples on $[0,1]\times X_1$, with $a_1$ invertible
outside $[0,1]\times K_1$, gives the triple $m(a_1,a_2)$ on $[0,1]\times X_1\times
X_2$, invertible outside $[0,1]\times K_1\times K_2$, which is a homotopy of the
products.

\emph{(e) Direct sums and stabilization.} If $\mathbb E^1=\mathbb E\oplus\mathbb E'$
with $a_1=a\oplus a'$, then $(\mathbb E^1\otimes\mathbb E^2)_\pm=
(\mathbb E\otimes\mathbb E^2)_\pm\oplus(\mathbb E'\otimes\mathbb E^2)_\pm$ after
reordering summands, and every entry of \eqref{families:eq:clifford} preserves the
splitting, so $m(a\oplus a',a_2)=m(a,a_2)\oplus m(a',a_2)$. Hence
$(x+x')\boxtimes x_2=x\boxtimes x_2+x'\boxtimes x_2$. For the stabilizing triple
$(H,H,\Id)$, $m(\Id_H,a_2)$ is invertible everywhere, so $[(H,H,\Id)]\boxtimes x_2=0$
by \ref{families:K-degenerate}; therefore stabilization does not change the product.

Since the equivalence relation is generated by these moves, $x_1\boxtimes x_2$ is
well defined, and by (e) (and its analogue in the second factor) it is additive in
each variable; additivity implies $(-x)\boxtimes y=-(x\boxtimes y)$, so the product is
bilinear.

\emph{The flip.} Let $\sigma:X_1\times X_2\to X_2\times X_1$ be the flip and
$s:F\otimes E\to E\otimes F$ the swap. The triple $y\boxtimes x$ has bundles
$F_+E_+\oplus F_-E_-$ and $F_-E_+\oplus F_+E_-$ (tensor signs omitted) and the map
\begin{equation*}
\begin{pmatrix}a_2\otimes1&-1\otimes a_1^*\\1\otimes a_1&a_2^*\otimes1\end{pmatrix}.
\end{equation*}
Pulling back by $\sigma$ and applying $s$, the source becomes $E_+F_+\oplus
E_-F_-=(\mathbb E^1\otimes\mathbb E^2)_+$ in the same order, and the target becomes
$E_+F_-\oplus E_-F_+$, which is $(\mathbb E^1\otimes\mathbb E^2)_-$ with the two
summands interchanged. Reordering the rows, $\sigma^*(y\boxtimes x)$ is represented by
\begin{equation*}
m''=\begin{pmatrix}a_1\otimes1&1\otimes a_2^*\\1\otimes a_2&-a_1^*\otimes1\end{pmatrix}
=\begin{pmatrix}a_1\otimes1&-1\otimes a_2^*\\1\otimes a_2&a_1^*\otimes1\end{pmatrix}
\begin{pmatrix}1&0\\0&-1\end{pmatrix}=(a_1\#a_2)\,J .
\end{equation*}
With $J_t=\operatorname{diag}(1,e^{i\pi t})$, $(a_1\#a_2)J_t$ is invertible wherever
$a_1\#a_2$ is, so it is a homotopy of triples from $a_1\#a_2$ to $m''$ (equivalently,
$J$ is an automorphism of $(\mathbb E^1\otimes\mathbb E^2)_+$ and the two triples are
isomorphic). Hence $\sigma^*(y\boxtimes x)=x\boxtimes y$.

\subsection*{\Cref{families:ex:CP1}}
\emph{Idea: for a rank one projection $\pi=uu^*$ with $\abs u=1$ one has
$\tr(\pi(d\pi)^2)=du^*\wedge du$.}

Let $u$ be a smooth unit vector field with $\pi=uu^*$ (locally $u=v/\abs v$), and
$f=u^*du$, a scalar $1$-form with $du^*\,u=-f$ (differentiate $u^*u=1$). Then
$d\pi=du\,u^*+u\,du^*$ and
\begin{equation*}
d\pi\wedge d\pi=du\,(u^*du)\,u^*+du\wedge du^*+u\,(du^*\wedge du)\,u^*+u\,(du^*u)\wedge du^* .
\end{equation*}
Using $\tr(uu^*X)=u^*Xu$ for matrix-valued forms $X$, the four terms contribute
$f\wedge f$, $f\wedge(-f)$, $du^*\wedge du$ and $(-f)\wedge(-f)$; the scalar $1$-form
$f$ has $f\wedge f=0$. Hence
\begin{equation*}
\tr\bigl(\pi(d\pi)^2\bigr)=du^*\wedge du .
\end{equation*}
Now $v=(1,w)$, $\lambda=\abs v^2=1+\abs w^2$, $u=gv$ with $g=\lambda^{-1/2}$. Then
\begin{equation*}
du^*\wedge du=(g\,dv^*+v^*dg)\wedge(g\,dv+v\,dg)
=g^2\,d\bar w\wedge dw+g\,dg\wedge(\bar w\,dw-w\,d\bar w),
\end{equation*}
since $dv^*\wedge dv=d\bar w\wedge dw$, $v^*dv=\bar w\,dw$, $dv^*v=w\,d\bar w$ and
$dg\wedge dg=0$. With $g\,dg=\frac12d(\lambda^{-1})=-\frac12\lambda^{-2}(w\,d\bar w+\bar
w\,dw)$,
\begin{equation*}
g\,dg\wedge(\bar w\,dw-w\,d\bar w)=-\tfrac12\lambda^{-2}\cdot2\abs w^2\,d\bar w\wedge dw,
\end{equation*}
so $\tr(\pi(d\pi)^2)=(\lambda^{-1}-\abs w^2\lambda^{-2})\,d\bar w\wedge dw=
\lambda^{-2}\,d\bar w\wedge dw$. As $d\bar w\wedge dw=(dx-i\,dy)\wedge(dx+i\,dy)=
2i\,dx\wedge dy$,
\begin{equation*}
\frac{i}{2\pi}\tr\bigl(\pi(d\pi)^2\bigr)=\frac{i}{2\pi}\cdot\frac{2i}{(1+\abs w^2)^2}
dx\wedge dy=-\frac1\pi(1+\abs w^2)^{-2}dx\wedge dy .
\end{equation*}
(At $w=0$ this is the value found in \cref{families:ex:grassmann}; the formula was
also checked symbolically.) The affine chart covers $\C P^1$ except one point, and
with the complex orientation $dx\wedge dy$,
\begin{equation*}
\int_{\C P^1}\frac{i}{2\pi}\tr\bigl(\pi(d\pi)^2\bigr)=-\frac1\pi\int_{\R^2}
\frac{dx\,dy}{(1+x^2+y^2)^2}=-\frac1\pi\cdot2\pi\int_0^\infty\frac{r\,dr}{(1+r^2)^2}
=-\frac1\pi\cdot2\pi\cdot\frac12=-1 .
\end{equation*}

\subsection*{\Cref{families:ex:bott-chern}}
\emph{Idea: with trivial bundles and connections the curvature of the bent connection
is exact, and Stokes reduces everything to $\frac{i}{2\pi}\oint a^{-1}da$ along the
boundary of the square, where only the top edge contributes.}

\emph{The representative.} Let $f$ be as in the definition of $b_1$ and choose
$\epsilon>0$ with $f=0$ on $[0,\epsilon]$ and $f=2\pi$ on $[2\pi-\epsilon,2\pi]$. The class
$b_1$ is represented by $(\C,\C,a)$ with $a$ smooth on $T^*I$,
\begin{equation*}
a(\theta,\xi)=e^{-if(\theta)}\ (\xi\ge1),\qquad a=1\ (\xi\le-1),\qquad a=1\ \text{for }
\theta\notin(\epsilon,2\pi-\epsilon),
\end{equation*}
extended arbitrarily to the compact set $K=[\epsilon,2\pi-\epsilon]\times[-1,1]$; $a$ is
invertible outside $K$ (these are exactly the conditions making its extension by $1$
to $T^*\Sphere$ the modified symbol of $P_1$). Take $R\ge2$ and
$X=[\epsilon/2,2\pi-\epsilon/2]\times[-R,R]$, so $a$ is invertible near $\pa X$. (The
square has corners; Stokes' theorem holds on it, or one can round the corners inside
$\{a\text{ invertible}\}$ without changing the computation below.)

\emph{The Chern character.} Take $E_\pm=\C$ with $\nabla_\pm=d$, so $\omega_\pm=0$, and
$\gamma=a^{-1}da$ where $a$ is invertible. Choose $\rho\in\CIc(\operatorname{int}X)$
with $\rho=1$ near $K$. Then $\nabla^\rho=d+(1-\rho)\gamma$ by
\eqref{families:eq:nabla-rho}, and since $\gamma$ is a scalar $1$-form,
$\gamma\wedge\gamma=0$ and $\omega_{\nabla^\rho}=d((1-\rho)\gamma)$. Thus
\begin{equation*}
\lambda_\rho=e^{c\,d((1-\rho)\gamma)}-1=c\,d\bigl((1-\rho)\gamma\bigr),\qquad c=\frac
i{2\pi},
\end{equation*}
where $(1-\rho)\gamma$ is defined on all of $T^*I$ because $1-\rho=0$ near $K$. Since
$d\gamma=-\gamma\wedge\gamma=0$, $\lambda_\rho=-c\,d\rho\wedge\gamma$ is supported in
$\supp d\rho\subset\operatorname{int}X$. By
Stokes and $\rho=0$ on $\pa X$,
\begin{equation*}
\int_{T^*I}\Ch(b_1)=\int_X\lambda_\rho=c\int_{\pa X}a^{-1}da .
\end{equation*}
This is the formula of \cref{families:thm:fedosov} with $\varphi=1$: the interior term
$\tr e^{c\omega_+}-\tr e^{c\omega_-}$ vanishes, and $\Omega(s)=-s(1-s)\gamma^2=0$, so the
boundary term is $c\int_{\pa X}\gamma$.

\emph{The boundary integral.} For the orientation $d\theta\wedge d\xi$, $\pa X$ is
traversed counterclockwise in the $(\theta,\xi)$-plane. On the bottom edge
($\xi=-R$) and on the two vertical edges ($\theta=\epsilon/2$, $\theta=2\pi-\epsilon/2$)
$a=1$, so they contribute nothing. On the top edge $\xi=R$, $\theta$ decreases from
$2\pi-\epsilon/2$ to $\epsilon/2$ and $a^{-1}da=-if'(\theta)\,d\theta$, so
\begin{equation*}
\int_{\pa X}a^{-1}da=\int_{2\pi-\epsilon/2}^{\epsilon/2}-if'(\theta)\,d\theta
=-i\bigl(f(\epsilon/2)-f(2\pi-\epsilon/2)\bigr)=2\pi i .
\end{equation*}
Hence $\int_{T^*I}\Ch(b_1)=\frac{i}{2\pi}\cdot2\pi i=-1$.

\subsection*{\Cref{families:ex:odd-product}}
\emph{Idea: the hint's homotopy joins $\operatorname{diag}(aa',1)$ to
$\operatorname{diag}(a,a')$, and the odd Chern forms of a smooth homotopy differ by
an exact form.}

Let $J=\bigl(\begin{smallmatrix}0&1\\-1&0\end{smallmatrix}\bigr)$ and
$R_t=\bigl(\begin{smallmatrix}\cos(\pi t/2)&\sin(\pi t/2)\\-\sin(\pi t/2)&\cos(\pi t/2)
\end{smallmatrix}\bigr)$, both acting on $\C^N\oplus\C^N$ (entries are multiples of
$\Id_N$); $R_t$ is constant on $X$, $R_0=\Id$, $R_1=J$. Set
\begin{equation*}
h_t=\begin{pmatrix}a&0\\0&1\end{pmatrix}R_t\begin{pmatrix}a'&0\\0&1\end{pmatrix}R_t^{-1}
\in\CI([0,1]\times X;\GL(2N,\C)).
\end{equation*}
Then $h_0=\operatorname{diag}(aa',1)$, and since
$J\operatorname{diag}(a',1)J^{-1}=\begin{pmatrix}0&1\\-a'&0\end{pmatrix}
\begin{pmatrix}0&-1\\1&0\end{pmatrix}=\operatorname{diag}(1,a')$,
$h_1=\operatorname{diag}(a,a')$.

By the transgression formula \eqref{families:eq:odd-transgression}, with
$\theta_t=h_t^{-1}dh_t$ ($d$ in the $X$ variables) and $\varphi_t=h_t^{-1}\pa_th_t$,
\begin{equation*}
\tr(\theta_1^{2k+1})-\tr(\theta_0^{2k+1})=\int_0^1\pa_t\tr(\theta_t^{2k+1})\,dt
=d\Bigl((2k+1)\int_0^1\tr(\varphi_t\theta_t^{2k})\,dt\Bigr),
\end{equation*}
an exact form. (For completeness: $\pa_t\theta=d\varphi+\theta\varphi-\varphi\theta$,
$d\theta=-\theta^2$ gives $d(\theta^{2k})=0$, and by \eqref{families:eq:graded-trace}
$\pa_t\tr(\theta^{2k+1})=(2k+1)\tr((\pa_t\theta)\theta^{2k})=(2k+1)\tr(d\varphi\,
\theta^{2k})=(2k+1)d\tr(\varphi\theta^{2k})$.) For block diagonal maps $\theta$ is
block diagonal, so
\begin{equation*}
\tr(\theta_0^{2k+1})=\tr\bigl(((aa')^{-1}d(aa'))^{2k+1}\bigr),\qquad
\tr(\theta_1^{2k+1})=\tr\bigl((a^{-1}da)^{2k+1}\bigr)+\tr\bigl((a'^{-1}da')^{2k+1}\bigr).
\end{equation*}
Multiplying by $c_k$ and summing over $k$, $\Chodd(a)+\Chodd(a')-\Chodd(aa')$ is exact.

\subsection*{\Cref{families:ex:pullback-index}}
\emph{Idea: the pulled-back family has the same operators on the same fibres, the
pulled-back sections stabilize it, and the frames of the null bundle pull back to
frames.}

\emph{Setting.} Let $f:B'\to B$ be smooth, $B'$ compact, and
$M'=\{(m,b');\phi(m)=f(b')\}\xrightarrow{\phi'}B'$ the pulled-back fibration
(\cref{product:prop:pullback}), with $F:M'\to M$, $F(m,b')=m$, which maps the fibre
$M'_{b'}$ diffeomorphically onto $Z_{f(b')}$. Bundles are pulled back by $F$. For
$A\in\Psido0(M/B;E_+,E_-)$, $f^*A$ is defined by $(f^*A)_{b'}=A_{f(b')}$ under this
identification. Over a trivializing set $U\subset B$, $\phi^{-1}(U)\cong Z\times U$ and
$\phi'^{-1}(f^{-1}U)\cong Z\times f^{-1}(U)$, and $f^*A$ corresponds to the family
$A\circ f\in\CI(f^{-1}(U);\Psido0(Z))$, smooth as the composite of smooth maps; so
$f^*A\in\Psido0(M'/B')$ by \cref{product:lem:local-families}. Its principal symbol is
the pull-back of $\sigma_0(A)$ under the proper map $T^*(M'/B')\to T^*(M/B)$ induced by
$F$ (an isomorphism on each fibre), so $f^*A$ is elliptic and
$[\sigma_0(f^*A)]=f^*[\sigma_0(A)]$.

\emph{Stabilization.} Let $S$ be given by $v_1,\dots,v_N\in\CI(M;E_-)$ with
$A_b\oplus S_b$ surjective for all $b$, and let $f^*S$ be given by $v_j\circ F$. On the
fibre over $b'$, $(f^*A)_{b'}\oplus(f^*S)_{b'}=A_{f(b')}\oplus S_{f(b')}$, which is
surjective. Hence $f^*S$ stabilizes $f^*A$, and as subspaces of
$\CI(Z_{f(b')};E_+)\oplus\C^N$,
\begin{equation*}
\mathcal N(f^*A,f^*S)_{b'}=\mathcal N(A,S)_{f(b')}=(f^*\mathcal N(A,S))_{b'} .
\end{equation*}
This fibrewise identity is an isomorphism of smooth vector bundles: by the proof of
\cref{families:lem:stabilization}(ii), near any $b_0$ there are local frames $e_1,\dots,
e_r$ of $\mathcal N(A,S)$ over $U\ni b_0$, with $e_i\in\CI(\phi^{-1}(U);E_+)\oplus
\CI(U;\C^N)$. The pulled-back sections $e_i\circ F$ (and $e_i^{\C^N}\circ f$) are smooth
on $\phi'^{-1}(f^{-1}U)$ and restrict on each fibre to a basis of
$\mathcal N(f^*A,f^*S)_{b'}$, so they are local frames of $\mathcal N(f^*A,f^*S)$; they are
also, by definition, local frames of $f^*\mathcal N(A,S)$. In these frames the identity
map has the identity matrix. Therefore, by \cref{families:lem:stabilization}(iii) (any
stabilization may be used),
\begin{equation*}
\ind(f^*A)=[f^*\mathcal N(A,S)]-[\C^N]=f^*\bigl([\mathcal N(A,S)]-[\C^N]\bigr)=f^*\ind(A),
\end{equation*}
since $f^*:K^0(B)\to K^0(B')$ is a homomorphism with $f^*[\C^N]=[\C^N]$. With
\ref{families:F-symbol} and \cref{families:lem:stabilization}(iv) this says
$\inda(f^*x)=f^*\inda(x)$ for all $x\in\Kc(T^*(M/B))$.

\emph{Consequence.} Let $\tilde\phi:\tilde M\to M$ have fibres $Y_m$, $b\in\Kc(T^*(\tilde
M/M))$ with $\inda(b)=1$, and $m\in M$. Apply the above to $i_m:\{m\}\hookrightarrow M$.
The pulled-back fibration is $Y_m\to\{m\}$, and the induced map
$T^*Y_m\hookrightarrow T^*(\tilde M/M)$ is the inclusion of a closed submanifold, hence
proper, so the restriction $b|_m=i_m^*b\in\Kc(T^*Y_m)$ is defined
(\ref{families:K-pullback}). Writing $b=[\sigma_0(P)]$ by \ref{families:F-symbol},
$b|_m=[\sigma_0(P_m)]$ and
\begin{equation*}
\ind(P_m)=\inda(b|_m)=i_m^*\inda(b)=i_m^*[\C]=1\in K^0(\{m\})=\Z,
\end{equation*}
where the first equality is \cref{families:rem:stab-remarks}(a) (for a point, the
index bundle is the numerical index). So the restriction of $b$ to every fibre has
numerical index $1$.

\subsection*{\Cref{families:ex:winding}}
\emph{Idea: bundles over $\R^2$ are trivial, $\det$ turns the moves of $\Kc$ into
operations that preserve the winding number, and $b_1$ was computed in
\cref{families:ex:bott-chern}; then compare two homomorphisms $\Z\to\Z$.}

Orient $\R^2$ and, through an orientation preserving diffeomorphism, $T^*I$ with the
coordinates $(\theta,\xi)$, by $d\theta\wedge d\xi$; circles are traversed
counterclockwise. For a nowhere vanishing $g$ on a closed curve $\Gamma$ put
$w(g,\Gamma)=\frac1{2\pi i}\oint_\Gamma g^{-1}dg\in\Z$.

\emph{Definition of $W$.} Every vector bundle over $\R^2$ (and over $[0,1]\times\R^2$)
is trivial, by \cref{bott:lem:bundle-homotopy}(ii). So every triple is isomorphic to one
of the form $(\C^r,\C^{r'},a)$; as $a$ is invertible somewhere, $r=r'$. Let $a$ be
invertible outside the compact $K$, and set $W(a)=w(\det a,C_R)$ for a circle $C_R=\{\abs
z=R\}$ with $K\subset\{\abs z<R\}$. This is independent of $R$: for $R<R'$ the form
$(\det a)^{-1}d\det a$ is closed on the annulus between $C_R$ and $C_{R'}$ (a scalar
$1$-form $g^{-1}dg$ satisfies $d(g^{-1}dg)=-g^{-2}dg\wedge dg=0$), so Stokes gives equal
integrals. The same Stokes argument shows that $C_R$ may be replaced by the positively
oriented boundary of any compact domain with piecewise smooth boundary that contains $K$
in its interior.

\emph{Invariance under the moves.} (i) Bundle isomorphisms: in the trivializations they
are maps $g_\pm:\R^2\to\GL(r,\C)$, and $a'=g_-ag_+^{-1}$ outside a compact set. Since
$w(\cdot,C_R)$ is additive for products and $\det g_\pm$ extends to the disc bounded by
$C_R$ without zeros, $w(\det g_\pm,C_R)=0$ (Stokes on the disc), so $W(a')=W(a)$. This also
shows that $W$ does not depend on the trivializations chosen. (ii) Changing $a$ on a compact
set: take $R$ larger. (iii) Homotopies $a_t$ on $[0,1]\times\R^2$, invertible outside
$[0,1]\times K$: on $C_R$, $t\mapsto\det a_t|_{C_R}$ is a homotopy of loops in $\C^*$, and
$w$ is a homotopy invariant (it is continuous in $t$ and integer valued). (iv) Stabilization:
$\det(a\oplus\Id_H)=\det a$. Since the equivalence relation is generated by these moves, $W$
is well defined on $\Kc(\R^2)$; since $\det(a\oplus a')=\det a\det a'$, it is a
homomorphism.

\emph{$W(b_1)=1$.} With the representative of \cref{families:ex:bott-chern}, $W(b_1)$ is
the winding number of $a$ along the boundary of the square $X$ there, which contains $K$ in
its interior. The computation there gives $\oint_{\pa X}a^{-1}da=2\pi i$, so $W(b_1)=1$.

\emph{The index on the circle.} Take $N=1$ and $B$ a point in
\cref{families:prop:normalization}: $\beta_1:\Z=K^0(\text{pt})\to\Kc(T^*I)$,
$\beta_1(n)=b_1\boxtimes n=nb_1$ (by bilinearity and $b_1\boxtimes1=b_1$, the product
with the triple $(\C,0)$ on a point being $b_1$ itself), and
$\inda\circ(k_1)_*\circ\beta_1=\Id_\Z$. By \ref{families:BP}, $\beta_1$ is an
isomorphism, so $\Kc(T^*I)=\Z\,b_1$. Both $\inda\circ(k_1)_*$ and $W$ are homomorphisms
$\Kc(T^*I)\to\Z$ taking $b_1$ to $1$, hence they coincide:
\begin{equation*}
\inda\bigl((k_1)_*x\bigr)=W(x),\qquad x\in\Kc(T^*I).
\end{equation*}
That is, if $A\in\Psido0(\Sphere;\C^r)$ is elliptic and $[\sigma_0(A)]=(k_1)_*x$ (its symbol
is, up to the moves of $\Kc$, the identity near $\theta=0$), then $\ind A$ equals the
winding number of $\det$ of the symbol along a large curve around the support of $x$,
traversed counterclockwise in the $(\theta,\xi)$-plane.
